\begin{hongqiarticle}{
        title  = {化悲痛为力量},
        author = {齐永红},
        style  = normal,
        body-size = 11.5pt,
        body-line-spacing = 1.6,
        body-indent = 2em
    }
    
    当代最伟大的马克思主义者，我党我军我国各族人民的伟大领袖，国际无产
    阶级和被压迫民族被压迫人民的伟大导师毛泽东主席与世长辞了。全国亿万军民
    怀着深厚的无产阶级感情，极其沉痛地举行了广泛的悼念活动，表达对伟大领袖
    毛主席的无限崇敬，寄托对伟大领袖毛主席的深切哀思。
    
    举国上下的掉念活动，在亿万军民的心中激发起巨大的力量。誓把毛主席开
    创的无产阶级革命事业进行到底的坚强决心，正在化为实际行动。工人含着热泪
    坚守岗位，抓批邓，促生产，创造了新的高产纪录；贫下中农含着热泪表示，要
    把毛主席亲手树立的人民公社这面红旗举得更高，不断巩固农村的社会主义阵
    地；人民解放军指战员和广大民兵含着热泪宣誓，誓死保卫党中央，保卫伟大祖
    国的无产阶级专政，加强战备，随时准备歼灭一切敢于入侵之敌，一定要解放台
    湾；上山下乡知识青年含着热泪发出豪言壮语，扎根农村一辈子，坚定不移地走
    毛主席指引的同工农相结合的道路；革命的知识分子含着热泪表决心，努力改造
    世界观，把文艺、教育，卫生等上层建筑领域的社会主义革命进行到底。全国各
    条战线，都在进一步学习毛主席著作，进一步实践毛主席指示。这一切有力地表
    明：在毛泽东思想哺育下成长起来的、经过无产阶级文化大革命、批林批孔、批邓、
    反击右倾翻案风斗争锻炼的中国人民，在这万分悲痛的时刻，是十分坚强的。
    
    在沉痛掉念毛主席逝世的日子里，我们要坚决响应党中央的号召，化悲痛为
    力量，继承毛主席的遗志，把毛主席所开创的并毕生为之奋斗的无产阶级革命事
    业坚持下去。
    
    化悲痛为力量，我们一定要加强学习。大海航行靠舵手，万物生长靠太阳，千
    革命靠毛泽东思想。这是我国人民经过长期革命实践所深深领会到的一条真理。
    半个多世纪以来，毛主席把马克思列宁主义的普遍真理和革命具体实践相结合，
    在同党内外、国内外阶级敌人的长期斗争中，继承、捍卫和发展了马克思列宁主
    义，发展了马克思主义的哲学、政治经济学和科学社会主义，极大地丰富了马克
    思主义的理论宝库。特别是毛主席关于无产阶级专政下继续革命的理论，解决
    了巩固无产阶级专政，防止资本主义复辟这个国际共产主义运动的重大课题，给
    予我们无产阶级革命事业的贡献是不可估量的。毛主席的一切著作和指示，都是
    马克思主义的光辉文献，是革命人民最宝贵的精神财富，是取之不尽，用之不竭
    的力量源泉。毛主席虽然和我们永别了，但毛主席的教导是永存的，战无不胜的
    毛泽东思想是永存的。我们一定要“\yulu{认真看书学习，弄通马克思主义}”，联系革命斗
    争的实践，刻苦攻读马列著作和毛主席著作。马克思主义、列宁主义、毛泽东思想
    为亿万人民群众所掌握，就会产生巨大的物质力量。为了继承毛主席的遗志，在
    沉痛悼念毛主席逝世的日子里，我们每个革命干部、党员和革命人民都要加倍努
    力地学习，使毛主席的教导永远成为我们行动的指南。
    
    化悲痛为力量，我们一定要更坚定地为执行和捍卫毛主席的无产阶级革命路
    线而战斗。“\yulu{思想上政治上的路线正确与否是决定一切的}。”这是毛主席率领我国人
    民进行五十多年革命所反复证明了的一条真理。毛主席全面总结了国际共产主义
    运动正反两个方面的经验，为我们党制定了一条马克思列宁主义的路线。我们的一
    切胜利，都是毛主席的无产阶级革命路线的胜利。我们党的全部历史表明，毛主席
    的无产阶级革命路线是我党我军我国各族人民的生命线。毛主席虽然和我们永别
    了，但毛主席的无产阶级革命路线是永存的，它象光芒万丈的灯塔，永远照耀着
    我们胜利前进的航程。毛主席嘱咐我们：“\yulu{按既定方针办}”。我们要永远不忘
    毛主席的嘱附，坚决按毛主席的无产阶级革命路线和各项政策办。在沉痛悼念
    毛主席逝世的日子里，我们每个革命干部、党员和革命人民，都要牢牢地树立一
    个信念：只要我们认真贯彻执行毛主席的革命路线，就有无穷无尽的力量，我们
    的事业就必然胜利。不论在任何时候和任何情况下，我们都要自觉坚持毛主席的
    革命路线，誓死捍卫毛主席的革命路线，同妄图改变党的基本路线的党内走资派
    展开长期不懈的斗争，把毛主席长期领导的反修防修的伟大斗争坚持到底。
    
    化悲痛为力量，我们一定要加强工人阶级领导的工农联盟为基础的各族人民
    的大团结，坚决维护党的团结和统一。由于毛主席关于在马克思列宁主义基础上增
    强党的团结的思想深入人心，我们党内的机会主义头子多次搞分裂，都没有搞成。
    每经过一次大的路线斗争，都使我们党在毛主席革命路线的基础上团结得更加坚
    强有力。在当前，我们要做到党政军民学，东西南北中，更紧密地团结在党中央
    的周围，服从党中央领导，听从党中央指挥，同国内外阶级敌人的一切破坏活
    动、分裂活动和阴谋诡计进行坚决的斗争。
    
    化悲痛为力量，我们一定要继续深入开展批邓、反击右倾翻案风的斗争。邓
    小平反革命的修正主义路线，是刘少奇、林彪反革命的修正主义路线的继续，它
    的矛头是对着毛主席和毛主席的无产阶级革命路线的。批邓，就是捍卫毛主席的
    无产阶级革命路线。邓小平煽起的右倾翻案风，翻文化大革命的案，算文化大革命
    的帐，妄图全盘否定毛主席亲自发动和领导的无产阶级文化大革命。批邓，就是
    捍卫和巩固无产阶级文化大革命的胜利成果。邓小平授意炮制的《论总纲》、《汇报
    提纲》和《条例》这三株大毒草，集中反映了邓小平反革命的修正主义路线的极右
    实质，集中反映了邓小平一类党内资产阶级的反动世界观。批邓，就是要抓住这
    三株大毒草，充分发挥邓小平这个反面教员的作用，使广大干部、党员和群众从
    路线上分清是非，提高识别真假马克思主义的能力。我们一定要以毛主席关于反
    击右倾翻案风的一系列重要指示为指针，发扬“\yulu{宜将剩勇追穷寇}”的革命精神，肃
    清邓小平反革命的修正主义路线的流毒，保证毛主席的革命路线在各条战线、各
    个领域得到进一步的贯彻。我们一定要以批邓为动力，以阶级斗争为纲，抓革
    命，促生产，促工作，促战备，把各项社会主义事业推向前进。
    
    化悲痛为力量，就是要以伟大的领袖和导师毛主席为光辉榜样，努力做毛主席
    的好学生，在大风大浪中把自己锻炼成为坚强的无产阶级革命战士。毛主席的一
    生，是彻底革命的一生。他为了中国人民的解放事业，为了全世界被压迫民族被
    压迫人民的解放事业，为了人类最壮丽的共产主义事业，贡献出了自已毕生的精
    力，直到生命的最后一息。我们沉痛悼念毛主席的逝世，就要认真地学习毛主席
    光辉一生的伟大革命实践，做到象毛主席教导的那样，全心全意为中国人民和世
    界革命人民谋利益，为实现共产主义奋斗终生；做到象毛主席教导的那样，相信
    群众，依靠群众，尊重群众的首创精神；做到象毛主席教导的那样，热情支持社会主
    义新生事物，关心新生力量的成长，自觉限制资产阶级法权，巩固和发展无产阶
    级文化大革命的胜利成果，做到象毛主席教导的那样，彻底革命，永远革命，发
    扬“\yulu{要扫除一切害人虫，全无敌}”的英雄气概；做到象毛主席教导的那样，敢于反
    潮流，“\yulu{不管风吹浪打，胜似闲庭信步}”，在国内同刘少奇、林彪、邓小平所代表的
    党内外资产阶级作顽强的斗争，在国际上同超级大国的霸权主义进行坚决的斗
    争，同以苏修叛徒集团为中心的现代修正主义进行坚决的斗争；做到象毛主席教
    导的那样，用辩证唯物主义和历史唯物主义的宇宙观观察问题，解决问题，同一
    切唯心论和形而上学作斗争。雨露滋润禾苗壮，毛泽东思想的阳光雨露，必将培
    育出一代文一代的无产阶级革命战士。
    
    化悲痛为力量，这是我们共产党人对自已的伟大领抽和导师开创的革命事业
    充满信心的表现。一八八三年三月十四日伟大的革命导师马克思逝世以后，恩格
    斯为这样一颗最强有力的心脏停止了跳动而感到万分悲痛。但是，恩格斯并没有
    为马克思的逝世失去信心，而是化悲痛为力量，继承马克思的遗志，强调“\yulu{我们决
        不会因此丧失勇气}”，坚信“\yulu{最后的胜利依然是确定无疑的}”。（《致弗·阿·左尔格》，
    《马克思恩格斯全集》第35卷第460页）马克思，恩格斯逝世以后，经历了同第二国际
    修正主义的斗争，无产阶级革命事业不但没有中断，相反地在全世界范围内有了
    蓬勃兴旺的发展；马克思主义的明灯不但没有熄灭，相反地在全世界范围内闪耀
    着更加灿烂的光辉。中国无产阶级在十月革命的影响下，找到了马克思列宁主义，
    建立了中国共产党。毛主席把马克思列宁主义的普遍真理同中国革命的具体实践
    相结合，引导我们党在阶级斗争、两条路线斗争中不断发展壮大。在我们党刚刚
    诞生的时候，全党只有几十个成员。在毛主席的英明领导下，经过五十五年的斗
    争，现在已经成为拥有三千多万党员、领导八亿人口、斗争经验极为丰富的马克思
    列宁主义的党。我们党从小到大，从弱到强，从领导新民主主义革命到领导社
    会主义革命的发展过程，就是坚持用马克思主义、列宁主义、毛泽东思想教育全
    党全国人民的过程。特别是经过无产阶级文化大革命，马克思主义、列宁主义、
    毛泽东思想得到了空前普及，毛主席的革命路线日益深入人心，千百万无产阶级
    革命事业的接班人茁壮成长。我们党及其领导的事业是大有希望的。共产党人永
    远向前，毛主席开创的无产阶级革命事业必将赢得更大的胜利。
    
    “\yulu{前途是光明的，道路是曲折的。}”毛主席经常用这样的话教育我们，要准
    备经过弯弯曲曲的道路去争取光明的前途。无产阶级必定战胜资产阶级，社会主
    义必定战胜资本主义，马克思主义必定战胜修正主义，共产主义一定要实现，这
    是不可抗拒的历史发展的总趋势。在沉痛悼念伟大的领袖和导师毛主席逝世的
    日子里，我们一定要把巨大的悲痛化为巨大的力量，继承毛主席的遗志，沿着
    毛主席指引的道路前进。下定决心，不怕牺牲，排除万难，为最终实现共产主义
    而奋斗！
\end{hongqiarticle}